% Part: first-order-logic
% Chapter: model-theory
% Section: partial-iso

\documentclass[../../../include/open-logic-section]{subfiles}

\begin{document}

\olfileid[fr]{mod}{bas}{dlo}
\section{Ordres linéaires denses}

\begin{defn}
  Un \emph{ordre linéaire dense sans extrémités} est
  !!a{structure}~$\Struct{M}$ pour le !!{language}
  ne contenant qu'un !!{predicate} binaire~$<$,
  qui satisfait les énoncés suivants :
  \begin{enumerate}
  \item $\lforall[x][\lnot x < x]$;
  \item $\lforall[x][\lforall[y][\lforall[z][(x < y \lif (y < z \lif x
    <z ))]]]$;
  \item $\lforall[x][\lforall[y][(x< y \lor \eq[x][y] \lor y < x)]]$;
  \item $\lforall[x][\lexists[y][x < y]]$;
  \item $\lforall[x][\lexists[y][y < x]]$;
  \item $\lforall[x][\lforall[y][(x < y \lif \lexists[z][(x < z \land
        z < y))]]]$.
 \end{enumerate}
\end{defn}

\begin{thm}\ollabel{thm:cantorQ}
  Deux ordres linéaires denses sans extrémités
  !!{enumerable}s sont isomorphes.
\end{thm}

\begin{proof}
  Soient $\Struct{M_1}$ et $\Struct{M_2}$ deux ordres linéaires
  denses sans extrémités !!{enumerable}s, avec
  ${<_1} = \Assign{<}{M_1}$ et ${<_2} = \Assign{<}{M_2}$.
  Soit $\PIso{I}$ l'ensemble de tous les isomorphismes
  partiels entre ces deux structures. Il est non vide,
  puisque $\emptyset \in \PIso{I}$. Montrons qu'il
  possède la propriété de va-et-vient; alors
  $\Struct{M_1} \iso[p] \Struct{M_2}$, et le théorème
  découlera du \olref[pis]{thm:p-isom1}.

  Pour la propriété d'aller, prenons $p \in \PIso{I}$
  et supposons $p(a_i) = b_i$ pour $i = 1$, \dots,~$n$,
  avec, quitte à réindexer,
  $a_1 <_1 a_2 <_1 \cdots <_1 a_n$.
  Soit $a \in \Domain{M_1}$. Si $a$ appartient déjà
  au domaine de $p$, prenons $q=p$. Sinon, si
  $n=0$, choisissons n'importe quel
  $b\in\Domain{M_2}$. Pour $n>0$, choisissons
  $b \in \Domain{M_2}$ selon la position de $a$ :\footnote{Complément éditorial : la source n'explicite ni le cas où $a$ est déjà dans le domaine, ni celui de l'application vide. Nous traitons ces deux cas avant les trois positions d'intervalle.}
  \begin{enumerate}
  \item si $a <_1 a_1$, prenons $b$ tel que $b <_2 b_1$;
  \item si $a_n <_1 a$, prenons $b$ tel que $b_n <_2 b$;
  \item si $a_i <_1 a <_1 a_{i+1}$ pour un certain $i$,
    prenons $b$ tel que $b_i <_2 b <_2 b_{i+1}$.
  \end{enumerate}
  Un tel $b$ existe parce que $\Struct{M_2}$ est un
  ordre linéaire dense sans extrémités. Dans les cas
  où $a$ n'appartient pas au domaine de $p$, posons
  $q = p \cup \{ \langle a, b \rangle \}$.
  On obtient $q \in \PIso{I}$, prolongement voulu de
  $p$. Cela établit la propriété d'aller. La propriété
  de retour se démontre de même. Ainsi,
  $\Struct{M_1} \iso[p] \Struct{M_2}$ et, par le
  \olref[pis]{thm:p-isom1}, $\Struct{M_1} \iso
  \Struct{M_2}$.
\end{proof}

\begin{prob}
  Compléter la démonstration du
  \olref[mod][bas][dlo]{thm:cantorQ}
  en vérifiant que $\PIso{I}$ possède la propriété de retour.
\end{prob}

\begin{rem}
  Soit $\Struct{S}$ un ordre linéaire dense sans
  extrémités !!{enumerable} quelconque. Par le
  \olref{thm:cantorQ}, on a $\Struct{S} \iso
  \Struct{Q}$, où $\Struct{Q} = (\Rat, <)$
  est l'ordre linéaire dense !!{enumerable} dont le
  !!{domain} est l'ensemble $\Rat$ des nombres rationnels.
  Reprenons la !!{structure}~$\Struct{R} =
  (\Real, <)$ de la \olref[thm]{remark:R}. Nous avons
  vu qu'il existe !!a{structure}~$\Struct{S}$
  !!{enumerable} telle que $\Struct{R} \elemequiv
  \Struct{S}$. Or $\Struct{S}$ est un ordre linéaire
  dense sans extrémités !!{enumerable} : elle est donc
  isomorphe, et par suite élémentairement équivalente,
  à $\Struct{Q}$. Par transitivité de l'équivalence
  élémentaire, $\Struct{R} \elemequiv
  \Struct{Q}$. On aurait aussi pu le montrer
  directement en établissant $\Struct{R}
  \iso[p] \Struct{Q}$ par le même argument
  de va-et-vient.
\end{rem}
\end{document}
